% ============================================================================
% Documentación: Colonia Artificial de Abejas simple (ABC) para CARP
% Generado/actualizado: 2026-07-10. Fuente: metacarp/abejas_simple.py,
% scripts/run_abc_simple_automatico.py, _calibracion_2knob.
% ============================================================================
\section{Colonia Artificial de Abejas simple (ABC)}
\label{sec:mh-abc}

\subsection{Concepto algorítmico}
El ABC (Karaboga, 2005) modela el forrajeo de una colmena. Cada
\textbf{fuente de alimento} es una solución candidata y su ``néctar'' es la
calidad (inverso del objetivo penalizado). Cada ciclo tiene tres fases:

\begin{enumerate}
  \item \textbf{Empleadas}: cada una de las \texttt{num\_fuentes} abejas
    empleadas genera un vecino de \emph{su} fuente; si mejora, la reemplaza y
    reinicia su contador de intentos \texttt{trials[i]}; si no,
    \texttt{trials[i]{+}{+}}.
  \item \textbf{Observadoras}: eligen fuentes por ruleta proporcional al
    fitness (las mejores fuentes reciben más intentos de mejora) y aplican el
    mismo criterio greedy de las empleadas.
  \item \textbf{Scouts}: toda fuente con $\texttt{trials} \ge
    \texttt{limite\_abandono}$ se considera agotada y se reemplaza por una
    solución \emph{aleatoria pura} (barajado de tareas + empacado greedy por
    capacidad), fiel al espíritu exploratorio del algoritmo original.
\end{enumerate}

La versión ``simple'' del proyecto es la referencia canónica de Karaboga con
una única modernización: el sesgo inter/intra-ruta compartido
($P(\text{inter}) = \max(p_{inter}, 0.8)$ bajo violación de capacidad), para
que la reparación de capacidad sea consistente con SA/TS/RTS. La fase de
observadoras se evalúa en lote (vectorizada, GPU opcional).

\subsection{Parámetros}
\begin{table}[htbp]
\centering
\caption{Parámetros del ABC simple.}
\begin{tabular}{lll}
\toprule
Parámetro & Significado & Valor calibrado \\
\midrule
\texttt{num\_fuentes} & tamaño de la población de fuentes & 30 (grid) \\
\texttt{limite\_abandono} & intentos sin mejora antes del scout & 60 (knob 2) \\
$p_{inter}$ & prob.\ de operadores inter-ruta (factible) & 0.5 (grid) \\
\texttt{iteraciones} & ciclos ABC (presupuesto) & instance-aware \\
\texttt{max\_iter\_sin\_mejora} & parada por estancamiento & instance-aware \\
\bottomrule
\end{tabular}
\end{table}

Por defecto \texttt{num\_fuentes} $\sim 2\sqrt{n_{tareas}}$ (la población
crece con la raíz de la dimensión); la grilla lo exploró vía factores de
escala.

\subsection{Búsqueda de malla (grid search)}
\begin{itemize}
  \item \textbf{Grilla principal} (\texttt{run\_abc\_simple\_automatico.py}):
    \texttt{factor\_fuentes} $\in \{1.5, 2.0, 3.0, 4.0\}$ (4) $\times$
    \texttt{factor\_abandono} $\in \{0.25, 0.5, 0.75, 1.0\}$ (4) $\times$
    $p_{inter} \in \{0.4, 0.5, 0.6, 0.7, 0.8\}$ (5) $\times$
    \texttt{factor\_iter} $\in \{15, 20, 30\}$ (3)
    $= 240$ combinaciones $\times$ 23 instancias $\times$ 5 repeticiones
    $= 27{,}600$ corridas. Los factores multiplican los defaults
    instance-aware, de modo que la grilla es comparable entre instancias de
    distinto tamaño.
  \item \textbf{Segundo knob}: \texttt{limite\_abandono} explorado con el
    resto fijo; ganó $60$ con gap medio de $9.529\,\%$.
  \item \textbf{Configuración fija resultante}: \texttt{num\_fuentes} $= 30$,
    \texttt{limite\_abandono} $= 60$, $p_{inter} = 0.5$; corrida con los tres
    approaches a 5 repeticiones.
\end{itemize}

\paragraph{Resultado en las 23 pequeñas.} Mejor-por-instancia: gap medio
$1.32\,\%$, 13/23 BKS, con \texttt{pr\_aislado} ganando 16/23 celdas y
\texttt{binario\_capacidad} las 7 restantes.

\subsection{Transferencia a las instancias val/egl}
Se mantiene \texttt{num\_fuentes} $= 30$ (con presupuesto de evaluaciones
alineado, más población implica menos ciclos; 30 es un buen compromiso). El
\texttt{limite\_abandono} sí debe crecer con la dimensión --- en el ABC
canónico el límite es proporcional a $SN \cdot D$ --- pues con 60 fijo las
fuentes de instancias grandes se abandonarían antes de poder refinarse:
\[
  \texttt{limite\_abandono} = \max(60,\ n_{tareas}),
\]
con $p_{inter} = 0.5$ (val) / $0.6$ (egl). Presupuesto uniforme: $10^6$
evaluaciones ($\texttt{iteraciones} = \lfloor 10^6 / (2 \cdot
\texttt{num\_fuentes}) \rfloor = 16{,}666$, pues cada ciclo evalúa
$\approx 2\,SN$ vecinos entre empleadas y observadoras) o 300\,s de pared.
Approach: \texttt{pr\_aislado}.
